\section{From the cover band to the Cayley projection}
\label{spec:section}

Lemma~\ref{spec:model} identifies the model's separated upper space;
Theorem~\ref{sam:realization} supplies a voltage-cover approximation to its
compression.
We now transfer the band to that cover and then to the Cayley graph.
Disjoint assigned domains give the spectral comparison, while all indexed
copies will give an equivariant comparison of traces.

Retain the gap $\chi$, deleted fraction $\beta_\ell$, model $O(t)$, and
upper projection $E_*(t)$ from the preceding section. With
$\mathfrak w_i=2^{C_\bullet H_i}$, Proposition~\ref{nil:weighted} gives
\begin{equation}
 \mu_\ell:=\sup_{\substack{0<t<\delta\\i\in J(t)}}
       \mathfrak w_i\,\operatorname{tr}_{\mathrm{id}}p(t)=o(1).
 \label{spec:weighted-input}
\end{equation}
The trace here and throughout this section includes the ordinary,
\emph{unnormalized} sum over the $N$ matrix indices.

Theorem~\ref{sam:realization} supplies, for any fixed accuracy $\eta>0$ and
all sufficiently large $\ell$, an actual self-adjoint weighted adjacency
$T_P$ on the voltage cover, satisfying
Equation~\eqref{sam:approx}.  Thus
\begin{equation}
  \|T_P-\Pi O\Pi\|\leq\eta,
  \label{spec:approx-input}
\end{equation}
where $\Pi$ deletes the removed base vertices, and $T_P$ is extended by zero
on those vertices.  Both operators in Equation~\eqref{spec:approx-input}
have layer width one.  Every edge weight before normalization has norm at
most a fixed $C_\eta$, the normalized weight is determined by its augmented
label, and the finite positive alphabet $S$ satisfies $|S|=o(dD)$.  The cover
satisfies the girth hypothesis of Theorem~\ref{geo:main}.  That Theorem
provides the indexed embedded copies, their disjoint assigned domains, the
finite trimming balls, the incidence count, the remaining forest, and the
torsion-freeness of $\Gbase$.  In particular these are conclusions proved
before this section, rather than additional spectral assumptions.

\subsection{Compression and norm perturbation}

We record quantitative Hilbert-space estimates, including the projection
equivalences needed later for the trace.

\begin{lemma}[Compression of a separated upper space]
\label{spec:compression}
Let $A=A^*$ on a Hilbert space, with $\|A\|\leq2$, and let $E$ be a
reducing projection such that
\[
 \|(A-I)E\|\leq a,\qquad
 A|_{(1-E)\mathcal H}\leq(1-b)I,\qquad 0<b<1.
\]
Let $J$ be an orthogonal projection with
$\varepsilon:=\|(1-J)E\|\leq1/2$.  The operator $JE$ has closed range.
Writing $E^J$ for the projection onto that range, put, on $J\mathcal H$,
\begin{equation}
 B_J=E^J+(J-E^J)A(J-E^J).
 \label{spec:compression-model}
\end{equation}
Then
\begin{equation}
 \|JAJ-B_J\|\leq a+20\varepsilon,
 \qquad B_J|_{(J-E^J)\mathcal H}\leq(1-b)I,
 \qquad \|B_J\|\leq2.
 \label{spec:compression-error}
\end{equation}
There is a partial isometry $V$ with $V^*V=E$, $VV^*=E^J$, and
$\|V-E\|\leq2\varepsilon$.  If $A,E,J$ commute with a specified unitary
or projective unitary action, then so do all the projections and maps just
constructed.  In particular, in a setting with the identity trace of
Lemma~\ref{nil:trace}, $E$ and $E^J$ have equal identity traces.  These
statements include $E=0$.
\end{lemma}

\begin{proof}
In the corner $E\mathcal H$,
\[
       (1-\varepsilon^2)E\leq EJE\leq E.
\]
Thus $JE$ is bounded below on $E\mathcal H$, and
\[
        V=JE(EJE)^{-1/2}
\]
is the asserted partial isometry.  The inverse is taken in the $E$ corner.
For $\varepsilon\leq1/2$, functional calculus gives
$\|(EJE)^{-1/2}-E\|\leq2\varepsilon^2$, whence
$\|V-E\|\leq\varepsilon+2\varepsilon^2\leq2\varepsilon$.
Every vector in $(J-E^J)\mathcal H$ is orthogonal to $E\mathcal H$:
if $Jy=y$ and $y\perp JE\mathcal H$, then
$\langle y,Ex\rangle=\langle y,JEx\rangle=0$.  This proves the
complementary upper bound in Equation~\eqref{spec:compression-error}.

Transporting the $E^J$ corner by $V$ yields
\[
 \|V^*AV-E\|
 \leq a+2\|A\|\,\|V-E\|
 \leq a+8\varepsilon.
\]
Reduction by $E$ and the preceding orthogonality also give
\[
 \|(J-E^J)AV\|
 =\|(J-E^J)A(V-E)\|\leq4\varepsilon.
\]
The difference between $JAJ$ and Equation~\eqref{spec:compression-model}
has just this pair of adjoint off-diagonal blocks and the estimated upper
diagonal block.  The norm of the paired off-diagonal blocks is the norm
of either one.  Its total norm is at most $a+12\varepsilon$, and therefore
at most the stated $a+20\varepsilon$.  The orthogonal two-part form of
$B_J$ gives $\|B_J\|\leq2$.  Equivariance follows from the formulas and
functional calculus.  Equality of the two traces follows from the
adjoint-square identity for $V$ in Lemma~\ref{nil:trace}.
\end{proof}

\begin{lemma}[An elementary contour estimate]
\label{spec:perturbation}
Suppose $B=B^*$ has a reducing projection $F$ with $BF=F$, and its
complementary spectrum is contained in $[-2,1-b]$, where $0<b<1$.
Let $r_0>0$ satisfy $r_0\leq b/2$ and $r_0<1$.  If
$A=A^*$ and $\|A-B\|\leq u<r_0$, the circle
$\mathcal C_{r_0}=\{z:|z-1|=r_0\}$ lies in the resolvent sets of both
operators.  The enclosed spectral projection $F_A$ satisfies
\begin{equation}
 \|F_A-F\|\leq\frac{u}{r_0-u}.
 \label{spec:projection-distance}
\end{equation}
Its enclosed spectrum lies in $[1-u,1+u]$; all the other spectrum of $A$
is at most $1-b+u$.  If $u<r_0/2$, the projections $F_A$ and $F$ are
equivalent by a partial isometry belonging to the algebra of bounded
operators generated by them.  This equivalence preserves any common
equivariance and its identity trace.
\end{lemma}

\begin{proof}
On $\mathcal C_{r_0}$ one has
$\|(z-B)^{-1}\|\leq r_0^{-1}$.  The resolvent Neumann series gives
$\|(z-A)^{-1}\|\leq(r_0-u)^{-1}$.  The resolvent identity and integration
around the circle then give Equation~\eqref{spec:projection-distance}.
For a real number at distance greater than $u$ from $\operatorname{spec}B$,
the same Neumann argument proves it is outside $\operatorname{spec}A$.
This yields the asserted spectral intervals.

For completeness, if projections $F_1,F_2$ have $\|F_1-F_2\|<1$, the map
$F_2F_1$ is bounded below on $F_1\mathcal H$, since
$\|F_2x\|\geq(1-\|F_1-F_2\|)\|x\|$ there.  Its range is all of
$F_2\mathcal H$: a vector in that space orthogonal to its range would
satisfy $F_1x=0$ and hence $\|(F_2-F_1)x\|=\|x\|$.  Its polar
normalization
\[
 F_2F_1(F_1F_2F_1)^{-1/2}
\]
therefore has initial projection $F_1$ and final projection $F_2$.
Apply this with $F_1=F$, $F_2=F_A$.  The equivariance and trace conclusions
follow from these formulas and Lemma~\ref{nil:trace}.
\end{proof}

Set
\begin{equation}
       \alpha=\chi/1000,\qquad r_0=\chi/3,
       \qquad C_\chi=4/\chi.
 \label{spec:accuracy-constants}
\end{equation}
Choose the fixed approximation accuracy $\eta>0$ so small that
\begin{equation}
 \eta\leq\alpha/2,\qquad
 \theta:=C_\chi\eta\leq1/2,\qquad
        2^{C_\bullet}\theta^2\leq1/2.
 \label{spec:eta-choice}
\end{equation}
The other fixed parameters of the construction have already been chosen;
Theorem~\ref{sam:realization} permits this accuracy.  We subsequently take
$\ell$ sufficiently large that
\begin{equation}
       20\sqrt{\beta_\ell}\leq\alpha/2.
 \label{spec:deletion-choice}
\end{equation}
All inequalities below concern such $\ell$ and are uniform in the layers.

\begin{proposition}[The actual cover projection and its exact fiber trace]
\label{spec:cover-projection}
The operator $T_P$ has a separated upper spectral projection $E_P$,
obtained by integration around $\mathcal C_{r_0}$.  It satisfies
\begin{equation}
 \|T_P\|\leq2,\qquad
 \|(T_P-I)E_P\|\leq\alpha,\qquad
 T_P|_{(1-E_P)\mathcal H}\leq(1-\chi+\alpha)I.
 \label{spec:cover-gap}
\end{equation}
In the central direct integral, for almost every $t$,
\begin{equation}
 \operatorname{tr}_{\mathrm{id}}E_P(t)
  =\begin{cases}
        \operatorname{tr}_{\mathrm{id}}p(t),&0<t<\delta,\\
        0,&t\notin(0,\delta).
    \end{cases}
 \label{spec:exact-fiber-trace}
\end{equation}
The projection $E_P$ is nonzero and reduces to the connected components of
$P$.  Its diagonal at a lifted vertex depends only on the base vertex.
\end{proposition}

\begin{proof}
Apply Lemma~\ref{spec:compression} to $O(t),E_*(t),\Pi$ with $a=0$ and
$b=\chi$.  Extend its two-part comparison operator by zero on
$(1-\Pi)\mathcal H_t^0$.  The extended complementary part is still at
most $1-\chi$, because $1-\chi>0$.  Denote its upper projection by
$E_*^\Pi(t)$.  Equations~\eqref{spec:approx-input},
\eqref{spec:deletion-choice}, and \eqref{spec:eta-choice} show that $T_P(t)$
differs from this comparison operator by at most $\alpha$ for almost
every $t$.  The passage from the full regular norm in
Equation~\eqref{spec:approx-input} to this almost-everywhere statement is
the direct-integral norm formula from the fiber construction.  Lemma~\ref{spec:perturbation}
with $r_0=\chi/3$ gives a common contour, Equation~\eqref{spec:cover-gap},
and
\[
       \|E_P(t)-E_*^\Pi(t)\|
              \leq\frac{\alpha}{r_0-\alpha}<1.
\]
All these fields are measurable by the resolvent formulas, and their
uniform bounds define the corresponding bounded direct-integral
operators.  The norm bound also follows directly from
$\|\Pi O\Pi\|\leq1$ and $\eta<1$.

Both the compression equivalence and the close-projection equivalence
commute with the fiber right actions used to define the identity trace.
Lemma~\ref{nil:trace} and Lemma~\ref{spec:model} therefore imply the
equality of traces in Equation~\eqref{spec:exact-fiber-trace}, with no
approximation error.  Outside the window the model is zero and
$\|T_P(t)\|\leq\eta$ almost everywhere.  Consequently its contour
projection is exactly zero there.  This conclusion uses the small norm
of the Fourier-mask leakage supplied by Theorem~\ref{sam:realization}; it
does not require the finite Fourier masks to vanish exactly.

Since $\operatorname{tr}_{\mathrm{id}}p(t)>0$ on the interval of positive
measure $(0,\delta)$, Equation~\eqref{spec:exact-fiber-trace} makes $E_P$
nonzero.  Every component subspace reduces the actual adjacency $T_P$,
and hence its resolvent and $E_P$.  Left deck translations commute with
$T_P$, since its entries are right deck translations with the prescribed
matrix coefficients.  They therefore commute with $E_P$.  They send any
lift of a fixed base vertex to every other lift, proving the last
assertion, even when they permute different connected components.
\end{proof}

\subsection{Localization of identity trace}

For a surviving base vertex $v$, let
\begin{equation}
 e_v=\sum_{j=1}^N
  \big\langle E_P(\delta_{(v,1_F)}\otimes e_j),
                \delta_{(v,1_F)}\otimes e_j\big\rangle\geq0,
 \qquad
 \mathcal S=\sum_{i\in\mathbb Z}\mathfrak w_i
                 \sum_{v\in V(Y)\cap(A_i\cup B_i)}e_v.
 \label{spec:weighted-size}
\end{equation}
Here $e_j$ denotes the standard matrix basis; put $e_v=0$ at a deleted
base vertex when summing over $Y^0$.  These are regular identity
diagonals before the central Fourier decomposition.

The weights record the two counts in
Equations~\eqref{geo:ball-weighted-sum} and~\eqref{geo:lost-weighted-sum}.
A fixed base type may have many lifts inside one trimming ball, while
the sum of losses over copies carries a ball-size factor from
Cauchy--Schwarz. Controlling $\mathcal S$ will keep both forms of
overlap loss small.

\begin{proposition}[The weighted size tends to zero]
\label{spec:localization}
For the fixed choice of $\eta$ in Equation~\eqref{spec:eta-choice},
\begin{equation}
                \mathcal S\leq96\mu_\ell=o(1).
 \label{spec:S-small}
\end{equation}
In particular the bound contains no additional factor equal to a layer
cardinality or to the matrix dimension already included in $\mu_\ell$.
\end{proposition}

\begin{proof}
Write $A_0(t)=\Pi O(t)\Pi$ and $\Delta(t)=T_P(t)-A_0(t)$.  The comparison
operator in the proof of Proposition~\ref{spec:cover-projection} is
within $\alpha/2$ of $A_0(t)$.  Therefore on $\mathcal C_{r_0}$,
\begin{equation}
 R_0(z,t)=(z-A_0(t))^{-1},\qquad
       \|R_0(z,t)\|\leq(r_0-\alpha/2)^{-1}\leq C_\chi.
 \label{spec:model-resolvent}
\end{equation}
Let $P_i$ be the projection to both sides of layer $i$, including all
transverse and matrix coordinates.  The operator $A_0(t)$ is supported
on the active layers $J(t)$.  Thus its resolvent is $z^{-1}I$ on the
orthogonal sum of the other layers, and has no matrix entries between
that sum and the active layers.  Since $\Delta(t)$ has layer width one,
its resolvent expansion is the norm-convergent series
\begin{equation}
 (z-T_P(t))^{-1}
       =\sum_{n=0}^\infty R_0(z,t)(\Delta(t)R_0(z,t))^n.
 \label{spec:resolvent-series}
\end{equation}
If $k=\operatorname{dist}(i,J(t))$ and $n<k$, all layer paths contributing
after multiplication on the left by $P_i$ remain outside $J(t)$.  Each
resolvent factor along such a path equals $z^{-1}I$; equivalently,
by induction on the number of width-one factors,
\[
 P_iR_0(z,t)(\Delta(t)R_0(z,t))^n
                   =z^{-n-1}P_i\Delta(t)^n\qquad(n<k).
\]
Its contour integral is zero, because the circle excludes zero.  Integrate
the remaining terms of Equation~\eqref{spec:resolvent-series} and use
Equation~\eqref{spec:model-resolvent}.  On the window this gives
\begin{equation}
     \|P_iE_P(t)\|
       \leq\frac{r_0C_\chi}{1-\theta}\theta^k
       \leq4\theta^{\operatorname{dist}(i,J(t))}.
 \label{spec:layer-norm}
\end{equation}
This estimate is also valid when $k=0$; there is then no vanishing initial
segment to remove.

Let $\mathfrak t_t$ denote the identity trace, summed over all base and
matrix indices of the fiber.  It applies to the countable amplification
by Lemma~\ref{nil:trace}.  Since $P_i$ and $E_P(t)$ are equivariant,
the adjoint-square identity and positivity imply
\begin{align}
 \mathfrak t_t(P_iE_P(t)P_i)
   &=\mathfrak t_t(E_P(t)P_iE_P(t))\notag\\
   &\leq\|P_iE_P(t)\|^2\,\mathfrak t_t(E_P(t))\notag\\
   &\leq16\theta^{2\operatorname{dist}(i,J(t))}
                    \operatorname{tr}_{\mathrm{id}}p(t).
 \label{spec:layer-trace}
\end{align}
The middle inequality follows from the operator inequality
$0\leq E_P(t)P_iE_P(t)\leq\|P_iE_P(t)\|^2E_P(t)$.  Thus it is the
\emph{total trace of the projection}, rather than the dimension of the
layer, that occurs.

The elementary bound $|H_i-H_j|\leq|i-j|$ implies
$\mathfrak w_i\leq\mathfrak w_j2^{C_\bullet|i-j|}$.
Assign every integer to one of its closest points in the nonempty set
$J(t)$.  At most two integers at each positive distance can be assigned
to a given point.  Since $|J(t)|\leq2$, Equation~\eqref{spec:eta-choice}
gives
\begin{align*}
 \sum_i\mathfrak w_i\theta^{2\operatorname{dist}(i,J(t))}
 &\leq2\max_{j\in J(t)}\mathfrak w_j
          \left(1+2\sum_{k\geq1}
                       (2^{C_\bullet}\theta^2)^k\right)\\
 &\leq6\max_{j\in J(t)}\mathfrak w_j.
\end{align*}
Multiplying Equation~\eqref{spec:layer-trace} by $\mathfrak w_i$ and
summing proves a pointwise bound of $96\mu_\ell$ for the weighted fiber
trace.  The projection vanishes off the window by
Equation~\eqref{spec:exact-fiber-trace}.  Central Fourier decomposition
identifies $e_v$ with the integral of its fiber identity diagonal.
Tonelli's Theorem applies to the nonnegative diagonals and sums, so
integrating over the probability circle proves
Equation~\eqref{spec:S-small}.
\end{proof}

\subsection{An actual element of the matrix reduced group algebra}

Use the finite alphabet supplied by Theorem~\ref{sam:realization} to form
the group $\Gbase$ of Theorem~\ref{geo:main}.  For a positive label
$\sigma$, let $W_\sigma$ be its matrix weight and let $D_\sigma=D_{ik}$
be its label-determined degree normalizer.  Put
\begin{equation}
 w_\sigma=(dD_\sigma)^{-1/2}W_\sigma,\qquad
 (T\xi)(g)=\sum_{\sigma\in S}
       \big(w_\sigma\xi(g\sigma)+w_\sigma^*\xi(g\sigma^{-1})\big).
 \label{spec:Cayley-operator}
\end{equation}
This is a finite sum of bounded matrix coefficients times right regular
translations.  It is self-adjoint, and inversion of the group coordinate
identifies it with an element of $M_N(\mathbb C[\Gbase])$ in the left
regular convention.  In particular
$T\in M_N(C_r^*(\Gbase))$.  We use the right convention in the proof,
so that $T$ commutes with left translations of $\Gbase$.

For an indexed copy $Q$, let $T_Q$ be the adjacency using just its own
formal edges.  The embedding and labels identify it with the restriction
of $T_P$ to the source component.  Let $E_Q$ be its upper spectral
projection on $\mathcal C_{r_0}$.  Equation~\eqref{spec:cover-gap}
holds for $T_Q,E_Q$, and if an incidence $(Q,u)$ has base type $v$,
\begin{equation}
           \operatorname{Tr}_{\mathbb C^N}E_Q(u,u)=e_v.
 \label{spec:copy-diagonal}
\end{equation}
The equation follows from component reduction and deck equivariance in
Proposition~\ref{spec:cover-projection}; it does not require the connected
component to be invariant under every deck translation.

\begin{lemma}[Finite trimming costs little operator norm]
\label{spec:copy-cut}
Let $D_Q$ and $B(Q)$ be the assigned domain and trimming ball from
Theorem~\ref{geo:main}.  Then
\begin{equation}
             \|\mathbf1_{B(Q)}E_Q\|^2\leq\mathcal S.
 \label{spec:ball-HS}
\end{equation}
For $\mathcal S\leq1/4$, the compression of $T_Q$ to
$\ell^2(D_Q)\otimes\mathbb C^N$ is within
$\alpha+20\sqrt{\mathcal S}$ of an operator $C_Q$ which is $1$ on a
projection $F_Q$ equivalent to $E_Q$, has norm at most $2$, and is at
most $1-\chi+\alpha$ on the complementary space.  In particular
$F_Q\ne0$ whenever $E_Q\ne0$.
\end{lemma}

\begin{proof}
The ball is finite.  For the ordinary Hilbert--Schmidt norm on vertex
and matrix indices, the projection identity gives
\[
 \|\mathbf1_{B(Q)}E_Q\|^2
 \leq\|\mathbf1_{B(Q)}E_Q\|_{\mathrm{HS}}^2
 =\sum_{u\in B(Q)}\operatorname{Tr}_{\mathbb C^N}E_Q(u,u).
\]
A fixed base vertex $v$ can occur at multiple lifts in the ball.  Its
multiplicity is at most $|B(Q)|$, and whenever it occurs in layer $i$,
the ball bound in Theorem~\ref{geo:main} gives
$|B(Q)|\leq\mathfrak w_i$.  Grouping the nonnegative sum by base vertex,
using Equation~\eqref{spec:copy-diagonal}, bounds it by
$\sum_i\mathfrak w_i\sum_{v\in A_i\cup B_i}e_v=\mathcal S$.
This proves Equation~\eqref{spec:ball-HS} with repetitions included.

The missing vertices $V(Q)\setminus D_Q$ lie in $B(Q)$.  Apply
Lemma~\ref{spec:compression} with $A=T_Q$, $E=E_Q$,
$J=\mathbf1_{D_Q}$, $a=\alpha$, $b=\chi-\alpha$, and
$\varepsilon\leq\sqrt{\mathcal S}$.  Its operator $B_J$ is the required
$C_Q$ and its projection $E^J$ is $F_Q$.  The explicit partial isometry
in Lemma~\ref{spec:compression} proves the nonvanishing assertion.
\end{proof}

\begin{lemma}[Matrix adjacency on a forest]
\label{spec:forest-norm}
Suppose a countable forest has maximum degree at most $m_1$, and a
self-adjoint adjacency on it has matrix edge weights of norms at most
$w_1$.  Its norm is at most $2w_1\sqrt{m_1}$.
\end{lemma}

\begin{proof}
Choose one root in each tree and orient every edge from parent to child.
Let $B$ carry the contribution from the parent coordinate to the child
coordinate, with the corresponding matrix weight.  Every nonroot vertex
has only one parent.  Hence for a finitely supported vector $\xi$,
\[
 \|B\xi\|^2
   \leq w_1^2\sum_x\#\{\text{children of }x\}\,\|\xi(x)\|^2
   \leq m_1w_1^2\|\xi\|^2.
\]
The self-adjoint adjacency is $B+B^*$, which proves the assertion first
on finitely supported vectors and then by bounded extension.  Infinite
trees cause no extra term; each connected tree is rooted at an ordinary
vertex and its finite-distance parent relation is well-defined.
\end{proof}

\begin{proposition}[A nonzero upper band in the reduced representation]
\label{spec:global-gap}
For sufficiently large $\ell$, the operator in
Equation~\eqref{spec:Cayley-operator} satisfies
\begin{equation}
 \operatorname{spec}T
 \subset(-\infty,1-\chi+4\alpha]\,
             \cup[1-3\alpha,1+3\alpha].
 \label{spec:global-spectrum}
\end{equation}
The second spectral part is nonempty.  Its projection
\begin{equation}
          p_0=\mathbf1_{[1-\chi/4,\infty)}(T)
                    \in M_N(C_r^*(\Gbase))
 \label{spec:p0}
\end{equation}
is consequently nonzero.
\end{proposition}

\begin{proof}
Increase $\ell$ so that
\begin{equation}
 \mathcal S<\min\{1/4,1/1000\},\qquad
                    20\sqrt{\mathcal S}\leq\alpha.
 \label{spec:S-final-choice}
\end{equation}
The assigned domains are disjoint.  Form the bounded orthogonal direct
sum
\[
       C=\bigoplus_Q C_Q\oplus0,
       \qquad F_D=\bigoplus_Q F_Q\oplus0,
\]
where the extra zero acts on all unassigned vertices.  The uniform norm
bound in Lemma~\ref{spec:copy-cut} proves these are bounded operators on
the ambient $\ell^2(\Gbase)\otimes\mathbb C^N$.  The direct sum $C$ is
$1$ on $F_D$, has norm at most $2$, and is at most $1-\chi+\alpha$ on
$1-F_D$.  The zero summand satisfies this upper bound because
$1-\chi+\alpha>0$.

At least one component of $P$ has a nonzero upper projection, since
$E_P\ne0$.  Its label maps give indexed copies, and
Lemma~\ref{spec:copy-cut} preserves the nonzero upper space in every such
copy.  Thus $F_D\ne0$.  This argument uses the ordinary direct sum and
does not assert that the domain assignment is equivariant.

Let $T_D$ denote the direct sum of the compressed $T_Q$, extended by zero
on unassigned vertices.  Lemma~\ref{spec:copy-cut} gives
$\|T_D-C\|\leq2\alpha$.  On every formal edge internal to an assigned
copy domain, $T_D$ has exactly the same matrix coefficient as $T$,
because that coefficient is determined by the label.  Disjointness of
the domains and the formal-edge embeddings prevent double counting of
these internal edges.  By the forest conclusion of
Theorem~\ref{geo:main}, the remaining formal edges form a forest.
Consequently $T-T_D$ is precisely a matrix adjacency of the kind in
Lemma~\ref{spec:forest-norm}.  Its maximum degree is at most $2|S|$ and
its weight norm is at most
\[
       w_1\leq\frac{\sqrt2 C_\eta}{\sqrt{dD}}.
\]
It follows that
\begin{equation}
       f_\ell:=\|T-T_D\|
          \leq4C_\eta\sqrt{\frac{|S|}{dD}}=o(1).
 \label{spec:forest-error}
\end{equation}
Take $\ell$ still larger so that $f_\ell\leq\alpha$.  Then
$\|T-C\|\leq3\alpha$.  Lemma~\ref{spec:perturbation}, with
$b=\chi-\alpha$ and the same $r_0=\chi/3$, proves
Equation~\eqref{spec:global-spectrum}.  It also shows that the enclosed
projection of $T$ has distance at most
$3\alpha/(r_0-3\alpha)<1$ from the nonzero projection $F_D$, so it is
nonzero.  Since $\alpha=\chi/1000$, the number $1-\chi/4$ lies strictly
between the two intervals in Equation~\eqref{spec:global-spectrum}.

The represented reduced algebra is a unital $C^*$-subalgebra of the
bounded operators and contains $T$.  For clarity, its spectrum for this
self-adjoint element equals the operator spectrum: a resolvent away
from the real spectrum is obtained by continuous functional calculus,
or by uniformly approximating the reciprocal on that compact spectrum
with polynomials in $T$ and the identity.  The function that equals zero
on the lower interval and one on the upper interval is continuous on
the spectrum.  Continuous functional calculus therefore gives precisely
the spectral projection in Equation~\eqref{spec:p0} as an element of
$M_N(C_r^*(\Gbase))$.
\end{proof}

\subsection{Equivariant size of all local upper spaces}

The comparison with the assigned domains proves spectral separation, but
their ordering need not be equivariant. Thus the ordinary equivalence with
$F_D$ does not by itself compare canonical traces. For the trace estimate
we return to all indexed copies. Their upper spaces form an equivariant
space of small trace, and we will show that the common kernel of the
corresponding restrictions contains no nonzero vector of the global upper space.

Write $\mathcal H=\ell^2(\Gbase)\otimes\mathbb C^N$. For a left
$\Gbase$-equivariant positive operator $Z$ on
$\ell^2(\Gbase)\otimes\mathcal K$, with a fixed countable orthonormal
basis $(k_a)$ of $\mathcal K$, its identity trace is
\begin{equation}
 \operatorname{Tr}_{\Gbase} Z
       =\sum_a\langle Z(\delta_1\otimes k_a),
                              \delta_1\otimes k_a\rangle\in[0,\infty].
 \label{spec:equivariant-trace}
\end{equation}
The size of a closed invariant subspace is the identity trace of its
orthogonal projection. Lemma~\ref{nil:trace} proves basis independence,
the adjoint-square trace identity, and invariance under equivariant
partial isometries in these countably amplified spaces. In particular
the size of $p_0\mathcal H$ is
\[
 \tau_{\Gbase,N}(p_0)
    =\sum_{j=1}^N\tau_{\Gbase}((p_0)_{jj}),
\]
with no division by $N$.

Consider the Hilbert space and projection
\begin{equation}
 \mathcal L=\bigoplus_Q\big(\ell^2(V(Q))\otimes\mathbb C^N\big),
       \qquad\mathcal E=\bigoplus_Q E_Q.
 \label{spec:all-copies}
\end{equation}
Left translation acts simultaneously on $Q$ and its vertices. The
incidence bijection of Proposition~\ref{geo:incidences} identifies this
representation, including matrix indices, with
\begin{equation}
 \mathcal L\cong\ell^2(\Gbase)\otimes\ell^2(V(Y))\otimes\mathbb C^N.
 \label{spec:incidence-identification}
\end{equation}
More explicitly, $(Q,u)$ is sent to $(u,v)$ where $v$ is its base type.
For each $u\in\Gbase$ and $v\in V(Y)$ there is exactly one such incidence:
choose any lift of $v$, map its component with that lift sent to $u$,
and use the deck-alignment equivalence in the definition of indexed
copies. Different lift choices give the same incidence, and any
incidence arises this way. The injective label maps ensure that the
base type is unambiguous. This count is valid for disconnected covers
and for copies with nontrivial, possibly infinite stabilizers. The
action on the pairs $(Q,u)$ is free because it is free on their group
vertex $u$; freeness of the action on the set of copies is unnecessary.

The operator $\mathcal E$ is equivariant in this representation. Its
identity diagonal at the base type $v$ has matrix sum $e_v$, by
Equation~\eqref{spec:copy-diagonal}. Hence
\begin{equation}
 \operatorname{Tr}_{\Gbase}\mathcal E
     =\mathcal S_0:=\sum_{v\in V(Y)}e_v
     \leq\mathcal S<\infty.
 \label{spec:all-copies-size}
\end{equation}
Thus $\mathcal U:=\mathcal E\mathcal L$ is a closed invariant space of
finite identity trace, even though there are countably many copies and
base vertices.

\begin{lemma}[The total loss from earlier intersections]
\label{spec:lost-restrictions}
Put $I_Q=V(Q)\setminus D_Q\subset B(Q)$. For every $\xi\in\mathcal H$,
\begin{equation}
 \sum_Q\|E_Q(\xi|_{I_Q})\|^2\leq\mathcal S\|\xi\|^2,
 \label{spec:all-copy-cut-loss}
\end{equation}
where restrictions are extended by zero inside their copies.
\end{lemma}

\begin{proof}
Write $v(Q,u)=\operatorname{type}_Q(u)$ for the base type of an incidence,
and $i(v)$ for the layer of a base vertex $v$. For a vector
$z\in\mathbb C^N$ supported at $u\in V(Q)$, positivity of the diagonal
matrix and Equation~\eqref{spec:copy-diagonal} give
\[
 \|E_Q(\delta_u\otimes z)\|^2
       =\langle E_Q(u,u)z,z\rangle
       \leq e_{v(Q,u)}\|z\|^2.
\]
The set $I_Q$ is finite. Cauchy--Schwarz for its sum therefore gives
\begin{equation}
 \|E_Q(\xi|_{I_Q})\|^2
       \leq |B(Q)|\sum_{u\in I_Q}
                  e_{v(Q,u)}\,\|\xi(u)\|^2.
 \label{spec:finite-cut-CS}
\end{equation}
The right-hand side is zero for an empty cut. If $u\in I_Q$ has base
type in layer $i$, the ball bound gives $|B(Q)|\leq\mathfrak w_i$.
Rearranging nonnegative terms yields
\begin{align*}
 \sum_Q\|E_Q(\xi|_{I_Q})\|^2
 &\leq\sum_{u\in\Gbase}\|\xi(u)\|^2
       \sum_{Q:u\in I_Q}|B(Q)|e_{v(Q,u)}\\
 &\leq\sum_{u\in\Gbase}\|\xi(u)\|^2
       \sum_{Q:u\in V(Q)}\mathfrak w_{i(v(Q,u))}e_{v(Q,u)}\\
 &=\mathcal S\|\xi\|^2.
\end{align*}
The final equality is the all-copy incidence count at the fixed vertex
$u$. In particular the factor $|B(Q)|$ from Cauchy--Schwarz is paid by
the layer weight in $\mathcal S$; it is not discarded or assumed
uniformly bounded.
\end{proof}

Together with disjointness of the $D_Q$, this controls the restriction
map to all upper spaces.

\begin{lemma}[A bounded restriction operator and a size bound]
\label{spec:restriction}
The formula
\begin{equation}
    A\xi=\big(E_Q(\xi|_{V(Q)})\big)_Q
       \quad(\xi\in\mathcal H)
 \label{spec:analysis-operator}
\end{equation}
defines a bounded equivariant operator $A:\mathcal H\to\mathcal U$,
with $\|A\|\leq1+\sqrt{\mathcal S}$. Its kernel $K$ is a closed
invariant subspace, and every $\xi\in K$ satisfies
\begin{equation}
        E_Q(\xi|_{V(Q)})=0\quad\text{for every indexed }Q.
 \label{spec:kernel-restrictions}
\end{equation}
Moreover,
\begin{equation}
           \operatorname{Tr}_{\Gbase} P_{K^\perp}\leq\mathcal S_0.
 \label{spec:kernel-complement-size}
\end{equation}
\end{lemma}

\begin{proof}
Since the $D_Q$ are pairwise disjoint and each $E_Q$ is a contraction,
\[
 \sum_Q\|E_Q(\xi|_{D_Q})\|^2
       \leq\sum_Q\|\xi|_{D_Q}\|^2\leq\|\xi\|^2.
\]
Together with Lemma~\ref{spec:lost-restrictions}, this shows that both
terms on the right in
\[
 A\xi=\big(E_Q(\xi|_{D_Q})\big)_Q
             +\big(E_Q(\xi|_{I_Q})\big)_Q
\]
belong to $\mathcal L$. The triangle inequality there gives the
asserted norm bound. Each coordinate belongs to the range of $E_Q$,
so the image is in the closed space $\mathcal U$. Although the
decomposition used to prove boundedness depends on the ordering of
copies, the formula~\eqref{spec:analysis-operator} does not. Left
translation permutes the copies and intertwines their $E_Q$, so $A$
is equivariant. Its kernel is therefore closed and invariant, and
Equation~\eqref{spec:kernel-restrictions} follows directly from the
coordinate formula.

The polar partial isometry of $A$ has initial projection
$P_{K^\perp}$ and final projection onto $\overline{\operatorname{ran}A}$,
a subspace of $\mathcal U$. It is equivariant by the polar-decomposition
argument in Lemma~\ref{nil:trace}. The adjoint-square identity and
monotonicity of that trace, together with
Equation~\eqref{spec:all-copies-size}, prove
Equation~\eqref{spec:kernel-complement-size}.
\end{proof}

\begin{lemma}[The upper global space injects into the small complement]
\label{spec:kernel-form}
With the choices in Equations~\eqref{spec:S-final-choice} and
\eqref{spec:forest-error},
\begin{equation}
 \langle T\xi,\xi\rangle
       \leq(1-\chi+3\alpha)\|\xi\|^2
              \quad(\xi\in K).
 \label{spec:kernel-form-bound}
\end{equation}
Consequently the bounded equivariant map
$P_{K^\perp}:p_0\mathcal H\to K^\perp$ is injective, and
\begin{equation}
        0<\tau_{\Gbase,N}(p_0)
               \leq\mathcal S_0\leq\mathcal S<1/2.
 \label{spec:strict-trace}
\end{equation}
\end{lemma}

\begin{proof}
For $\xi\in K$, Equation~\eqref{spec:kernel-restrictions} gives
\[
       E_Q(\xi|_{D_Q})=-E_Q(\xi|_{I_Q}).
\]
Lemma~\ref{spec:lost-restrictions} therefore implies
\begin{equation}
 \sum_Q\|E_Q(\xi|_{D_Q})\|^2\leq\mathcal S\|\xi\|^2.
 \label{spec:domain-upper-loss}
\end{equation}
The local spectral estimates in Equation~\eqref{spec:cover-gap} imply,
for every vector $y$ on $Q$,
\[
 \langle T_Qy,y\rangle
   \leq(1-\chi+\alpha)\|y\|^2+\chi\|E_Qy\|^2.
\]
Apply this to the zero-extended restrictions to $D_Q$, sum over the
disjoint domains, and use Equation~\eqref{spec:domain-upper-loss}.
The coefficient $1-\chi+\alpha$ is positive, so the unassigned part can
be included in its norm bound. Equation~\eqref{spec:forest-error} gives
\[
 \langle T\xi,\xi\rangle
 \leq(1-\chi+\alpha+\chi\mathcal S+f_\ell)\|\xi\|^2
 \leq(1-\chi+3\alpha)\|\xi\|^2,
\]
because $\mathcal S\leq1/1000$ and $f_\ell\leq\alpha$. This proves
Equation~\eqref{spec:kernel-form-bound} for every $\xi\in K$ directly;
no approximation by finitely supported kernel vectors is used.

On $p_0\mathcal H$, Equation~\eqref{spec:global-spectrum} gives
$\langle T\xi,\xi\rangle\geq(1-3\alpha)\|\xi\|^2$.
Since $1-3\alpha>1-\chi+3\alpha$, this upper space intersects $K$
only at zero. The map $P_{K^\perp}$ restricted to that space is
therefore injective. Both projections are equivariant, so its polar
decomposition identifies $p_0\mathcal H$ with a closed invariant
subspace of $K^\perp$. Equations~\eqref{spec:kernel-complement-size}
and~\eqref{spec:all-copies-size} prove the non-strict upper bounds in
Equation~\eqref{spec:strict-trace}.

Finally a nonzero equivariant projection on
$\ell^2(\Gbase)\otimes\mathbb C^N$ has strictly positive identity
trace. If all its identity diagonal entries were zero, their values as
squared projection-column norms would make all columns at the identity
zero; equivariance would make every group translate of those columns
zero, and hence the projection itself zero. Proposition~\ref{spec:global-gap}
proved $p_0\ne0$, so the lower bound is strict. The final upper bound
follows from Equation~\eqref{spec:S-final-choice}.
\end{proof}

\begin{theorem}[The matrix projection supplied by the construction]
\label{spec:matrix-projection}
For the fixed choices made above and some sufficiently large integer
$\ell$, the construction defines a finitely generated torsion-free
group $\Gbase$, a positive integer $N$, and the self-adjoint finite
group-algebra matrix $T$ in Equation~\eqref{spec:Cayley-operator}.  Its
spectrum in the reduced representation has the separation in
Equation~\eqref{spec:global-spectrum}.  The resulting projection $p_0$
in Equation~\eqref{spec:p0} belongs to $M_N(C_r^*(\Gbase))$ and obeys
\[
          0<\sum_{j=1}^N\tau_{\Gbase}((p_0)_{jj})<\frac12.
\]
\end{theorem}

\begin{proof}
The alphabet is finite by Theorem~\ref{sam:realization}, so the group
presentation in Theorem~\ref{geo:main} has finitely many generators;
Theorem~\ref{geo:main} proves torsion-freeness.  Equation~\eqref{spec:Cayley-operator}
is a finite sum in the actual reduced regular representation.
Proposition~\ref{spec:global-gap} proves its spectral separation and
the reduced-algebra membership and nonvanishing of $p_0$.
Lemma~\ref{spec:kernel-form} supplies the stated unnormalized trace
inequality.  The order of choices is consistent: the geometry and
transverse parameters are fixed first; Equation~\eqref{spec:eta-choice}
then fixes a single accuracy; the uniform vanishing of
$\beta_\ell$, $\mu_\ell$, and $|S|/(dD)$ permits all subsequent displayed
inequalities by one sufficiently large $\ell$.  No limit group or
limit projection is taken in this final step.
\end{proof}
